\subsection{সম্পর্কিত সারণিতে শিক্ষার্থীর মোট ফি}
\textbf{নমুনা।} একটি কলেজে নিচের দুই টেবিল আছে। এক শিক্ষার্থী বিভিন্ন মাসে ফি দেয়; প্রতিটি রসিদ আলাদা।
\begin{flushleft}
\small
\begin{tabular}{|l|l|l|}\hline
\multicolumn{3}{|l|}{Student}\\ \hline
StudentID & Name & ClassName\\ \hline
1 & Ayesha & XI\\\hline
2 & Nabil & XI\\\hline
3 & Rupa & XII\\ \hline
\end{tabular}
\qquad
\begin{tabular}{|l|l|l|}\hline
\multicolumn{3}{|l|}{Payment}\\ \hline
PaymentID & StudentID & Amount\\ \hline
501 & 1 & 600\\\hline
502 & 1 & 400\\\hline
503 & 2 & 700\\ \hline
\end{tabular}
\end{flushleft}

\begin{flushleft}
\small
\fbox{\parbox{0.88\linewidth}{\raggedright
Student.StudentID (PK) $\xrightarrow{\;1:N\;}$ Payment.StudentID (FK)}}
\end{flushleft}
\textbf{মোট জমা কমপক্ষে ৭০০ এমন শিক্ষার্থী নির্বাচন:}
\begin{verbatim}
SELECT s.StudentID, s.Name, SUM(p.Amount) AS TotalPaid
FROM Student AS s
INNER JOIN Payment AS p ON s.StudentID = p.StudentID
GROUP BY s.StudentID, s.Name
HAVING SUM(p.Amount) >= 700
ORDER BY TotalPaid DESC, s.StudentID ASC;
\end{verbatim}
\begin{flushleft}
\small
\begin{tabular}{|l|l|l|}\hline
StudentID & Name & TotalPaid\\ \hline
1 & Ayesha & 1000\\\hline
2 & Nabil & 700\\ \hline
\end{tabular}
\end{flushleft}
Ayesha-র মোট ৬০০+৪০০=১০০০, Nabil-এর ৭০০। Rupa-র Payment নেই, তাই INNER JOIN-এর ফলেও নেই। একই নামে দুই শিক্ষার্থী থাকতে পারে বলে শুধু Name দিয়ে গ্রুপ না করে StudentID-ও দেওয়া হয়েছে। “কমপক্ষে” মানে \texttt{>=}; শুধু \texttt{>} দিলে ৭০০ বাদ যেত।

